\lesson{3}{Oct 6 2021 Wen (09:29:49)}{The US Constitution and Your State}{Unit 2}

\subsubsection*{Road Trip Stop One- The U.S. Constitution}

Let's look at the basics of the Constitution.

\paragraph{Who} Delegates from  $12$ of the $13$ colonies met in Philadelphia. Only \textbf{Rhode Island} did not send a delegation. Many delegates were lawyers. \textbf{James Madison} ultimately proposed that the convention's goal should be to form a new government. The delegates debated how to form the best centralized government without restricting the rights of the states or individuals. The convention ultimately came down to compromises. Some famous members of the \textbf{Constitutional Convention} are:

\begin{itemize}
    \item George Washington
    \item Benjamin Franklin
    \item George Mason
    \item Roger Sherman
    \item Edmond Randolph
    \item James Madison
    \item Alexander Hamilton
\end{itemize}

Only $39$ of the $55$ delegates signed the U.S. Constitution.

\paragraph{What} The U.S. Constitution is the founding document of our government. The Constitution is a \textbf{"living document"} that provides the structure of our democracy. Our Constitution can be broken down into three major parts:

\begin{itemize}
    \item The Preamble
    \item The Articles
    \item The Amendments
\end{itemize}

We get the three branches of government from the Constitution.
\textbf{The Preamble} is the introduction to the Constitution that outlines the goals of the Constitution. The preamble is a $52$ word sentence that sets forth the goals of the U.S. government.

\paragraph{When} In the summer of $1787$, The members of the \textbf{Constitutional Convention} took $100$ days to frame our new government. The delegates debated and compromised on how our new government should function. One of the major debates centered on power. The power of the central government vs. the power of states' government.

\paragraph{Where} The delegates convened in \textbf{Philadelphia, Pennsylvania} at \textbf{Independence Hall}. This meeting of delegates is called the \textbf{Constitutional Convention} or \textbf{Philadelphia Convention}. This is also where the \textbf{Declaration of Independence} was signed.

\paragraph{Why} Their purpose was to establish a legal system of government that corrected the weakness of the \textbf{Articles of Confederation} and create a new centralized government. This was our second attempt to create our government.

\subsubsection*{Road Trip Stop Two – State Governments}

Now, let's dig a little bit deeper into the state and the local government.
\textbf{Article IV} of the Constitution could also be called the \textbf{"States Article"}. The $10$th amendment to the Constitution defines the powers of the states.

The $13$ original colonies maintained sovereignty by operating a self-government. However, during the time of the revolution, the independent states needed to work together to form a union, which would eventually create the \textbf{United States}. After the founding fathers saw the \textbf{Articles of Confederation} fail miserably, they needed to create a centralized government but still maintain the rights of each state.

\begin{definition}[Federalism]
    Federalism is the division of power between the \textbf{Federal} and \textbf{State} governments. The Constitution outlines the powers and limits of the \textbf{National Government}, solidifies the relationship between the levels of government and ensures that the people get their rights.
\end{definition}

\paragraph{Delegated to the National Government}

\begin{itemize}
    \item Organize and control trade between states and with other nations
    \item Coin and print money
    \item Make agreements with foreign nations
    \item Establish post offices and interstate highways
    \item Raise and support military forces
    \item Declare war and make peace
    \item Govern U.S. territories and admit new states and create rules and laws for people who wish to move into the nation and how they may become citizens
\end{itemize}

\paragraph{Concurrent or Shared Powers}

\begin{itemize}
    \item Raise and collect taxes
    \item Borrow money
    \item Establish courts
    \item Grant permissions to start banks
    \item Enforce laws and punish lawbreakers
    \item Provide health care services and other assistance to the people
\end{itemize}

\paragraph{Reserved to the State Governments}

\begin{itemize}
    \item Create local government
    \item Set up businesses
    \item Build and support public schools
    \item Organize and control trade within the state
    \item Conduct elections
    \item Determine rules for voting such as living in the district where the person votes
    \item Make marriage laws
    \item Determine the requirements for professional workers, like doctors and teachers
\end{itemize}

\textbf{State Governments} are also governed by the Constitution. In general, the \textbf{State Constitutions} are much longer than the U.S. Constitution. Each state:

\begin{itemize}
    \item Creates
    \item Amend
    \item Fund laws
\end{itemize}

However, the U.S. Constitution is the \textbf{"Supreme Law of the Land"}. 

\begin{enumerate}
    \item State governments receive federal funds for some services like transportation (driver's license), infrastructure (roads, bridges), and health care.
    \item State governments have three branches like the federal government (separation of power).
    \item The governor of the state is the head of the executive branch. The governor can veto state laws and is in charge of many departments (Executive Branch).
    \item Most states have a bicameral legislature (only Nebraska has a one house). State legislators write and pass laws for the citizens. The term limits and requirements vary by state (Legislative Branch).
    \item Citizens in a geographic region elect state legislators.
    \item Many state governments allow direct democracy, or citizens to directly add initiative to the ballot. Direct democracy also includes recall election and referendum.
    \item The judicial branch can vary by state. For example, the levels or courts or how judges get on the bench (elected or appointed).
    \item State governments deal with issues of public policy that impact citizens. They include education, death penalty, welfare, marriage laws, and voting standards.
\end{enumerate}

\subsubsection*{Road Trip Stop Three – Local Governments}

\begin{enumerate}
    \item Local governments get their power from the state government and are in place to carry out state initiatives.
    \item Local governments provide many services to its citizens such as police, fire departments, schools, parks, and supervision of elections.
    \item Local governments can be formed in cities, counties, or towns, etc.
    \item Local governments employ a lot of citizens, and are also voted into office by the citizens.
    \item Local governments usually have a council, a mayor/chairperson, and community representatives that focus on the local community and the issues they face.
    \item Local elections are held by postal elections. What is your 9–digit zip code?
\end{enumerate}

\newpage
